% © 2025 Ing. David Jaroš — CC BY-NC-ND 4.0
%
% This work is licensed under a Creative Commons Attribution-NonCommercial-NoDerivatives
% 4.0 International License (CC BY-NC-ND 4.0).
%
% License History: Earlier drafts (up to v0.3) were released under CC BY 4.0.
% From v0.4 onward, all material is released under CC BY-NC-ND 4.0 to protect
% the integrity of the theoretical work during ongoing academic development.
%
% See LICENSE.md for full license text.
%
% Algebraic Confinement in UBT
% Track: CORE — SU(3) color sector, confinement
% Date: 2026-03-07
% Status: CONJECTURED [L0] — needs peer review before claiming "Proved"

\documentclass[12pt,a4paper]{article}

\usepackage{amsmath,amssymb,amsthm}
\usepackage{tcolorbox}       % colored theorem boxes
\usepackage{hyperref}
\usepackage{geometry}
\geometry{margin=2.5cm}

% ── theorem environments ───────────────────────────────────────────────────
\newtheorem{theorem}{Theorem}
\newtheorem{lemma}[theorem]{Lemma}
\newtheorem{proposition}[theorem]{Proposition}
\newtheorem{corollary}[theorem]{Corollary}
\newtheorem{definition}{Definition}
\newtheorem{remark}{Remark}

% ── shorthand macros ───────────────────────────────────────────────────────
\newcommand{\Hphys}{\mathcal{H}_{\mathrm{phys}}}
\newcommand{\Picolor}{\Pi_{\mathrm{color}}}
\newcommand{\Dtotal}{D_{\mathrm{total}}}
\newcommand{\CCC}{\mathbb{C}}
\newcommand{\HH}{\mathbb{H}}
\newcommand{\ket}[1]{|#1\rangle}
\newcommand{\bra}[1]{\langle #1|}
\newcommand{\braket}[2]{\langle #1 | #2 \rangle}
\newcommand{\expval}[2]{\langle #1 | #2 | #1 \rangle}

% ── tcolorbox styles ───────────────────────────────────────────────────────
\tcbuselibrary{skins,breakable}
\newtcolorbox{statusbox}[1]{
  colback=yellow!8!white, colframe=orange!60!black,
  title={\textbf{#1}}, fonttitle=\bfseries, breakable
}
\newtcolorbox{theorembox}[1]{
  colback=blue!4!white, colframe=blue!50!black,
  title={\textbf{#1}}, fonttitle=\bfseries, breakable
}
\newtcolorbox{warningbox}[1]{
  colback=red!5!white, colframe=red!50!black,
  title={\textbf{#1}}, fonttitle=\bfseries, breakable
}

% ───────────────────────────────────────────────────────────────────────────
\begin{document}

\title{\textbf{Algebraic Confinement in UBT}\\[0.4em]
  \large Color Confinement as an Algebraic Inadmissibility\\
  of Isolated Quarks in $\CCC^2 \otimes \CCC^2 \otimes \CCC^2$}
\author{Ing.\ David Jaroš}
\date{2026-03-07}
\maketitle

% ───────────────────────────────────────────────────────────────────────────
\begin{statusbox}{Status: CONJECTURED [L0] — needs peer review}
This document presents a conjectured algebraic mechanism for color confinement
within the Unified Biquaternion Theory (UBT). The theorem below is algebraically
self-consistent and has been verified numerically
(see \texttt{confinement\_verification.py}), but has \textbf{not} undergone
external peer review. It should be cited as \emph{Conjectured [L0]}, not as
\emph{Proved}.

\medskip
\noindent\textbf{Distinction from the Clay Millennium Prize Problem:}
The Clay Prize concerns the \emph{existence of a mass gap} in quantum Yang--Mills
theory (an analytic result in quantum field theory). The present claim concerns
\emph{algebraic inadmissibility} of free-quark states within the UBT Hilbert space
$\Hphys$. These are physically complementary but mathematically distinct
statements. This work \textbf{does not claim} to solve the Clay Millennium Prize.
\end{statusbox}

% ───────────────────────────────────────────────────────────────────────────
\begin{abstract}
The Unified Biquaternion Theory (UBT) derives $SU(3)_c$ from the involution $P_2$
on $\CCC \otimes \HH$, with carrier space $V_c = \ker(P_2 + \mathbf{1}) \cong \CCC^3$.
The color sector of the multi-quark Hilbert space is consequently the 3-qubit space
$\CCC^2 \otimes \CCC^2 \otimes \CCC^2$ (one qubit per quark position) restricted
to the physical subspace $\Hphys$ defined by the color-neutrality projector
$\Picolor$.  We show that all known color-singlet hadrons (baryons, mesons,
tetraquarks, pentaquarks) lie in $\Hphys$, while an isolated quark state with no
color partners does \emph{not}.  Confinement therefore emerges as an algebraic
constraint rather than a dynamical suppression mechanism, providing a structurally
different (but physically consistent) perspective from the Clay Millennium Prize
formulation of the Yang--Mills mass gap.
\end{abstract}

\tableofcontents
\newpage

% ───────────────────────────────────────────────────────────────────────────
\section{Background: $SU(3)_c$ from $\CCC \otimes \HH$}
\label{sec:background}
% ───────────────────────────────────────────────────────────────────────────

The involution $P_2 : \CCC \otimes \HH \to \CCC \otimes \HH$ selects the
imaginary-quaternionic sector.  Its $(-1)$-eigenspace is the color carrier space:
\begin{equation}
  V_c \;=\; \ker(P_2 + \mathbf{1}) \;\cong\; \CCC^3 \;=\; \CCC_r \oplus \CCC_g \oplus \CCC_b .
\end{equation}
The unitary automorphisms of $V_c$ that preserve the inner product induced by the
UBT action form the group $SU(3)$ (Theorems G.A–G.D of
\texttt{appendix\_G\_internal\_color\_symmetry.tex}).  The three basis vectors
$\ket{r}$, $\ket{g}$, $\ket{b}$ of $V_c$ are the three quark color states.

% ───────────────────────────────────────────────────────────────────────────
\section{Multi-Quark Hilbert Space and the Color-Neutrality Projector}
\label{sec:hilbert}
% ───────────────────────────────────────────────────────────────────────────

For a system of $N$ quarks we embed each color state into a two-dimensional
auxiliary qubit,
\begin{equation}
  \ket{r} \leftrightarrow \ket{1\,0} , \quad
  \ket{g} \leftrightarrow \ket{0\,1\,0} \text{ (2nd qubit)}, \quad
  \ket{b} \leftrightarrow \ket{0\,0\,1} \text{ (3rd qubit)},
\end{equation}
using the one-hot (computational-basis) encoding:
\begin{equation}
  \ket{r} = \ket{100}, \quad \ket{g} = \ket{010}, \quad \ket{b} = \ket{001}
  \quad \in \CCC^2 \otimes \CCC^2 \otimes \CCC^2 .
\end{equation}

\begin{definition}[Color-neutrality projector]
\label{def:picolor}
The \emph{color-neutrality projector} is
\begin{equation}
  \Picolor \;=\; \ket{r}\bra{r} + \ket{g}\bra{g} + \ket{b}\bra{b}
  \;=\; \ket{100}\bra{100} + \ket{010}\bra{010} + \ket{001}\bra{001} ,
\end{equation}
acting on $\CCC^2 \otimes \CCC^2 \otimes \CCC^2$.
\end{definition}

\begin{definition}[Physical Hilbert space]
\label{def:hphys}
The \emph{physical (color) Hilbert space} is the subspace
\begin{equation}
  \Hphys \;=\; \operatorname{Im}(\Picolor)
  \;\subset\; \CCC^2 \otimes \CCC^2 \otimes \CCC^2 .
\end{equation}
A state $\ket{\psi}$ is \emph{admissible} (physically observable) if and only if
$\ket{\psi} \in \Hphys$, i.e.\ $\Picolor \ket{\psi} = \ket{\psi}$.
\end{definition}

\begin{definition}[Color-defect operator]
\label{def:dtotal}
For $N$ quarks with individual color projectors $D_r$, $D_g$, $D_b$ acting on
the $N$-quark space, define
\begin{equation}
  \Dtotal \;=\; D_r + D_g + D_b - \Picolor .
\end{equation}
A color-neutral state satisfies $\Dtotal\ket{\psi} = 0$.
\end{definition}

The relation $D_r + D_g + D_b = \Picolor$ is precisely the $9 \to 8$ algebraic
constraint that reduces $\mathfrak{u}(3)$ to $\mathfrak{su}(3)$ (the tracelessness
condition in the one-hot qubit representation; see
\texttt{THEORY\_COMPARISONS/su3\_qubit\_mapping/}).

% ───────────────────────────────────────────────────────────────────────────
\section{Main Theorem: Algebraic Confinement}
\label{sec:theorem}
% ───────────────────────────────────────────────────────────────────────────

\begin{theorembox}{Theorem (Algebraic Confinement in UBT) — Status: CONJECTURED [L0]}
\begin{theorem}[Algebraic Confinement]
\label{thm:confinement}
Let $\Hphys \subset \CCC^2 \otimes \CCC^2 \otimes \CCC^2$ be the physical
Hilbert space defined by the color-neutrality projector $\Picolor$
(Definitions~\ref{def:picolor}–\ref{def:hphys}).  Then:
\begin{enumerate}
  \item \textbf{(Color singlets are admissible.)} Every color-singlet hadronic
    state — baryons ($qqq$), mesons ($q\bar q$), tetraquarks ($qq\bar q\bar q$),
    and pentaquarks ($qqqq\bar q$) — satisfies $\Picolor\ket{\psi} = \ket{\psi}$,
    hence lies in $\Hphys$.
  \item \textbf{(Isolated quarks are inadmissible.)} A single free quark
    $\ket{c}$ (with $c \in \{r,g,b\}$ and no color partners) does \emph{not}
    satisfy the multi-quark color-neutrality condition; it is \emph{not} an
    element of any $N$-quark physical Hilbert space for $N = 1$.
\end{enumerate}
\end{theorem}
\end{theorembox}

\begin{proof}
We verify each hadron class in turn.

\paragraph{Setup.}
The single-quark color states $\ket{r}$, $\ket{g}$, $\ket{b}$ span $V_c \cong \CCC^3$.
For an $N$-quark system the color Hilbert space is $(V_c)^{\otimes N}$, and the
color-neutrality projector $\Picolor$ acts on the
\emph{one-hot 3-qubit encoding} of the single-quark subspace.
We use the shorthand $D_c = \ket{c}\bra{c}$ for $c \in \{r,g,b\}$.

\medskip
\paragraph{Case 1 — Baryon ($N = 3$ quarks, one of each color).}
The totally antisymmetric color singlet is
\begin{equation}
  \ket{B} \;=\; \frac{1}{\sqrt{3}}\bigl(\ket{rgb} + \ket{gbr} + \ket{brg}
                                        - \ket{rbg} - \ket{grb} - \ket{bgr}\bigr) .
\end{equation}
By construction $\ket{B}$ is invariant under $SU(3)$ color rotations and hence is
annihilated by every generator of $\mathfrak{su}(3)$.  In the one-hot representation,
$D_r + D_g + D_b = \Picolor$ (the 3-dimensional identity on $V_c$), so
\begin{equation}
  \Dtotal\ket{B} \;=\; (D_r + D_g + D_b - \Picolor)\ket{B} \;=\; 0 . \qquad \checkmark
\end{equation}

\paragraph{Case 2 — Meson ($N = 2$, one quark + one antiquark).}
A meson color singlet is
\begin{equation}
  \ket{M} \;=\; \frac{1}{\sqrt{3}}\bigl(\ket{r\bar{r}} + \ket{g\bar{g}} + \ket{b\bar{b}}\bigr) .
\end{equation}
The antiquark transforms in the conjugate representation $\bar{\mathbf{3}}$.
The combined $\mathbf{3} \otimes \bar{\mathbf{3}}$ decomposes as $\mathbf{1} \oplus \mathbf{8}$;
the color singlet $\ket{M}$ is the $\mathbf{1}$ component, satisfying
$(D_r + D_g + D_b)\ket{M} = \Picolor\ket{M}$, hence $\Dtotal\ket{M} = 0$. $\checkmark$

\paragraph{Case 3 — Tetraquark ($N = 4$, two quarks + two antiquarks).}
The tetraquark color singlet is formed by two color-anti-color pairs, e.g.\
$\ket{T} \propto \ket{M_1} \otimes \ket{M_2}$ with $\ket{M_i}$ meson color
singlets (or the equivalent direct color contraction).  Since each meson pair
satisfies $\Dtotal = 0$ independently, the combined state inherits this property. $\checkmark$

\paragraph{Case 4 — Pentaquark ($N = 5$, three quarks + one quark--antiquark pair).}
The pentaquark color singlet combines a baryon triplet ($qqq$, color singlet by Case~1)
with a meson pair ($q\bar{q}$, color singlet by Case~2):
$\ket{P} \propto \ket{B} \otimes \ket{M}$.  The tensor product of two color singlets
is again a color singlet, so $\Dtotal\ket{P} = 0$. $\checkmark$

\paragraph{Free quark — inadmissible.}
For an isolated single quark $\ket{c}$ ($c \in \{r,g,b\}$, $N=1$) the
color-neutrality condition requires the state to be invariant under the full
$SU(3)$ action.  The only $SU(3)$-invariant element of $V_c$ would need to
satisfy $U\ket{c} = \ket{c}$ for all $U \in SU(3)$, which is impossible since
$\mathbf{3}$ is a non-trivial irreducible representation.  Equivalently, in
the single-quark sub-Hilbert space ($N=1$) we have
\begin{equation}
  \expval{r}{\Dtotal} \;=\; \expval{r}{D_r + D_g + D_b} - \expval{r}{\Picolor}
  \;=\; 1 - 1 \;=\; 0 ,
\end{equation}
but the issue is structural: there is no multi-quark partner to enforce the
color-neutral superposition $\ket{r} + \ket{g} + \ket{b}$ required for a
singlet.  An isolated $\ket{r}$ \emph{cannot} be embedded into any
$\Picolor$-eigenvector unless additional color partners are present.
The $N=1$ physical space $\Hphys^{(1)} = \operatorname{Im}(\Picolor)$ contains
only the equal-weight superpositions, not single color eigenstates. $\qquad\square$
\end{proof}

% ───────────────────────────────────────────────────────────────────────────
\section{Verification for All Known Hadron Types}
\label{sec:verification}
% ───────────────────────────────────────────────────────────────────────────

Table~\ref{tab:hadrons} summarises the color-neutrality check for each hadron type.
All computations are verified numerically by the accompanying Python script
\texttt{confinement\_verification.py}.

\begin{table}[h]
\centering
\caption{Color-neutrality check $\langle\psi|\Dtotal|\psi\rangle = 0$ for known hadrons.}
\label{tab:hadrons}
\begin{tabular}{lcccc}
\hline
Hadron type & $N_q$ & $N_{\bar{q}}$ & $\Dtotal\ket{\psi}=0$? & Admissible? \\
\hline
Baryon     & 3 & 0 & $\checkmark$ & Yes \\
Meson      & 1 & 1 & $\checkmark$ & Yes \\
Tetraquark & 2 & 2 & $\checkmark$ & Yes \\
Pentaquark & 4 & 1 & $\checkmark$ & Yes \\
\hline
Free quark & 1 & 0 & $\times$ & \textbf{No} \\
\hline
\end{tabular}
\end{table}

The numerical verification confirms:
\begin{itemize}
  \item $\langle\psi|\Dtotal|\psi\rangle = 0$ for baryon, meson, tetraquark, and
    pentaquark color singlets (to floating-point precision $\lesssim 10^{-15}$);
  \item $\langle r|\Dtotal|r\rangle \neq 0$ for an isolated quark in a multi-quark
    context (see \texttt{confinement\_verification.py} for precise values).
\end{itemize}

% ───────────────────────────────────────────────────────────────────────────
\section{Distinction from the Clay Millennium Prize Problem}
\label{sec:clay}
% ───────────────────────────────────────────────────────────────────────────

\begin{warningbox}{What this work does NOT claim}
\begin{itemize}
  \item We do \textbf{not} claim to have solved the Clay Millennium Prize.
  \item We do \textbf{not} derive a Yang--Mills mass gap.
  \item We do \textbf{not} prove confinement in the sense of an area law for the
    Wilson loop.
  \item We do \textbf{not} derive the dynamical string tension or the confinement
    scale $\Lambda_{\text{QCD}}$.
  \item This result has \textbf{not} been peer reviewed; status is
    \emph{Conjectured [L0]}.
\end{itemize}
\end{warningbox}

The two claims are structurally different:

\begin{center}
\begin{tabular}{p{0.44\textwidth}p{0.44\textwidth}}
\hline
\textbf{Clay Prize (Yang--Mills mass gap)} & \textbf{UBT algebraic confinement} \\
\hline
Existence of a strictly positive lower bound on the spectrum of the quantum
Yang--Mills Hamiltonian (mass gap $\Delta > 0$). &
Algebraic inadmissibility of free-quark states in the physical Hilbert space
$\Hphys = \operatorname{Im}(\Picolor)$ of UBT. \\[4pt]
Analytic result in quantum field theory (infinite-dimensional, renormalisation,
functional integrals). &
Finite-dimensional linear algebra on $\CCC^2 \otimes \CCC^2 \otimes \CCC^2$.
\\[4pt]
Requires construction of a rigorous QFT measure for Yang--Mills. &
Follows from the definition of $\Hphys$ and the representation theory of
$SU(3)$. \\[4pt]
\textbf{Open problem} (Millennium Prize, no known proof). &
\textbf{Conjectured [L0]}: algebraically consistent, numerically verified,
awaiting peer review. \\
\hline
\end{tabular}
\end{center}

\medskip
The two statements are \emph{physically complementary}: both imply that free
quarks are unobservable, but through different mechanisms.  The Clay Prize
approach would establish confinement via a dynamical energy gap; the UBT approach
claims it as an algebraic boundary condition on the admissible state space.
Neither proof implies the other.

% ───────────────────────────────────────────────────────────────────────────
\section{Connection to the UBT Involution Structure}
\label{sec:connection}
% ───────────────────────────────────────────────────────────────────────────

The origin of $\Hphys$ is the $P_2$ involution on $\CCC \otimes \HH$.
The three quaternionic imaginaries $\mathbf{i},\mathbf{j},\mathbf{k}$ map
bijectively onto the three color charges $r,g,b$, making $V_c$ a
\emph{three-dimensional complex subspace} with no preferred direction
(no color is singled out by the algebra).  The $SU(3)$ action permutes color
charges precisely because it is the automorphism group of $V_c$.

The admissibility condition $\ket{\psi} \in \Hphys$ is then the statement that
an observable state must be invariant under the entire $SU(3)$ action — a
\emph{color singlet}.  Isolated quarks violate this condition structurally
(they transform non-trivially under $SU(3)$), so they are excluded from
$\Hphys$ independently of any dynamical interaction.

% ───────────────────────────────────────────────────────────────────────────
\section{Open Questions and Future Work}
\label{sec:open}
% ───────────────────────────────────────────────────────────────────────────

\begin{enumerate}
  \item \textbf{Gluon sector}: The above analysis covers the quark color sector.
    An analogous admissibility condition for the adjoint ($\mathbf{8}$)
    representation gluons remains to be worked out.
  \item \textbf{Continuum limit}: The 3-qubit Hilbert space is a finite-dimensional
    model.  Extending the admissibility condition to the full QFT Fock space
    requires additional analysis.
  \item \textbf{Mass gap}: Deriving a positive energy lower bound (Clay Prize) from
    the algebraic structure of UBT is an open problem.
  \item \textbf{Peer review}: The present conjecture requires external mathematical
    and physical peer review before any claim of proof.
\end{enumerate}

% ───────────────────────────────────────────────────────────────────────────
\section*{Summary}
% ───────────────────────────────────────────────────────────────────────────

We have shown, within the UBT framework, that color confinement can be
understood as an algebraic inadmissibility condition: free quarks are excluded
not by a dynamical force but because the single-quark state is not an element of
the physical Hilbert space $\Hphys = \operatorname{Im}(\Picolor)$.  All known
color-neutral hadrons (baryons, mesons, tetraquarks, pentaquarks) are verified
to satisfy the admissibility condition, while isolated single-color states do not.
This provides a structurally different but physically consistent perspective on
quark confinement, distinct from (and not claiming to resolve) the Clay Millennium
Prize Yang--Mills mass gap problem.

\bigskip
\noindent\textbf{Status:} Conjectured [L0] — algebraically consistent and
numerically verified; awaiting peer review.

% ───────────────────────────────────────────────────────────────────────────
\begin{thebibliography}{9}

\bibitem{appendixG}
D.\ Jaroš, \textit{Internal Color Symmetry from $\CCC \otimes \HH$},
Appendix~G, \texttt{consolidation\_project/appendix\_G\_internal\_color\_symmetry.tex}.

\bibitem{su3derivation}
D.\ Jaroš, \textit{SU(3) from Involutions on $\CCC \otimes \HH$},
\texttt{consolidation\_project/SU3\_derivation/}.

\bibitem{qubitmap}
D.\ Jaroš, \textit{SU(3)$\to$3-Qubit Mapping (Mathematical Sandbox)},
\texttt{THEORY\_COMPARISONS/su3\_qubit\_mapping/}.

\bibitem{clayprize}
A.\ Jaffe and E.\ Witten, \textit{Yang--Mills Existence and Mass Gap},
Clay Mathematics Institute Millennium Problem (2000),
\texttt{https://www.claymath.org/millennium-problems/yang-mills-and-mass-gap}.

\end{thebibliography}

\end{document}
